\subsection{Counting Techniques}
\hrulefill

\subsubsection*{Equally Likely Outcomes}
If all outcomes in a sample space are equally likely for any event A, then
\begin{equation*}
    P(A) = \frac{\text{Number of outcomes in A}}{\text{Number of outcomes in S}} = \frac{n(A)}{n(S)}
\end{equation*}
where $n(A)$ is the number of outcomes in event A and $n(S)$ is the number of outcomes in sample space S.

\subsubsection*{Combinatorics}
Combinatorics is the branch of mathematics dealing with combinations of objects in specific sets under certain constraints. It provides tools for counting and arranging objects.
Some basic combinatorial concepts include:
\begin{itemize}
    \item \textbf{Permutations:} The number of ways to arrange a set of objects in a specific order.
    \item \textbf{Combinations:} The number of ways to select a subset of objects from a larger set, without regard to the order of selection.
\end{itemize}

\paragraph*{Counting Subsets}
Let $\Omega = { a, b, c, d, e, f, g, h , i, j }$. How would we find all possible subsets of $\Omega$?
\begin{align*}
    Num = 2^{len(\Omega)} = 2^{10} = 1024
\end{align*}
This is because each element can either be included or excluded (2 options) from a subset, leading to $2^{n}$ possible subsets for a set with $n$ elements.

\pagebreak

\subsubsection*{Permutations and Combinations}
\paragraph*{Permutations}
A permutation is an arrangement of objects in a specific order. The number of permutations of 
$k$ objects chosen from $n$ distinct objects is given by the following formula:
\[
    P_{k,n} = \frac{n!}{(n-k)!} = n \times (n-1) \times (n-2) \times ... \times (n-k+1)
\]
Where $P_{k,n}$ is the number of permutations of $n$ objects taken $k$ at a time, $n!$ is the factorial of $n$, and $(n-k)!$ is the factorial of $(n-k)$.

\paragraph*{Example:} Suppose we want to rank 5 candidates Amy, Bai, Catherine, David, and Eduardo in preference for a position.
How many different ways can we rank (1.) all of the candidates and (2.) only 2 of the candidates?
\begin{align*}
    1. \quad P_{5,5} &= \frac{5!}{5-5!} = \frac{5!}{1} = 120 \\
    2. \quad P_{5,2} &= \frac{5!}{5-2!} = \frac{5!}{3!} = 20
\end{align*}

\paragraph*{Combinations}
A combination is a selection of objects without regard to the order of selection. The number of combinations
of $k$ objects chosen from $n$ distinct objects is given by the following formula:
\[
    C_{k,n} = \frac{n!}{k!(n-k)!} = \frac{P_{k,n}}{k!}
\]
Where $C_{k,n}$ is the number of combinations of $n$ objects taken $k$ at a time, $n!$ is the factorial of 
$n$, $k!$ is the factorial of $k$, and $(n-k)!$ is the factorial of $(n-k)$.

\paragraph*{Example:}
How many hands of 3 cards can be drawn from a standard deck of 52 playing cards?
\[
    C_{3,52} = \binom{52}{3} = \frac{52!}{3!(52-3)!} = \frac{52!}{3!49!} = 22100
\]

\paragraph*{Example:}
A playlist contains 100 songs. 10 of these songs are by The Strokes. We can find the probability of the first 
song by The Strokes being the 5th song played by finding the ratio of the number of favorable outcomes to the total number of outcomes.

\begin{align*}
    P(\text{1st Strokes on 5th}) &= \frac{\text{Number of favorable outcomes}}{\text{Total number of outcomes}} \\
    &= \frac{P_{4,90} \cdot P_{1,10}}{P_{5, 100}} = \frac{\frac{90!}{(90-4)!} \cdot 10}{\frac{100!}{(100-5)!}}\\
    &= .0679
\end{align*}

Another way of thinking about this would be that there are $\binom{100}{10}$ different ways to choose the 
10 Strokes songs. Now if we choose 9 of the last 95 songs to be Strokes songs, there are $\binom{95}{9}$ ways to do this,
 leaving 4 non-Strokes songs and 1 Strokes song in the first 5 songs. The probability is then:
\[
    P(\text{1st Strokes on 5th}) = \frac{\binom{95}{9}}{\binom{100}{10}} = .0679
\]

\paragraph*{Example:}
A university warehouse has recieved 25 printers, of which 10 are laser printers and 15 are inkjet printers.
If 6 printers are randomly selected for testing, what is the probability that exacty 3 of those selected are
laser printers?\\

Let $D_3$ be the event that exactly 3 of the selected printers are laser printers. 
\begin{align*}
    P(D_3) &= \frac{\text{Number of favorable outcomes}}{\text{Total number of outcomes}} \\
    &= \frac{N(D_3)}{N} = \frac{\binom{15}{3} \cdot \binom{10}{3}}{\binom{25}{6}} = .3083
\end{align*}